\section{Data Lake and Datasets}
\label{sec:data-lake}

\subsection{Data-Lake Scope}

The experimental data lake was constructed around New York City
urban-mobility analytics for January 2024. The core catalog contains four
physical datasets: Yellow Taxi trip records, Green Taxi trip records, Taxi
Zones, and hourly weather observations.

The choice deliberately combines datasets with different analytical roles.
Yellow and Green Taxi contain high-volume transactional trip records, Taxi
Zones provides geographic reference data, and the weather source provides a
temporal environmental dimension. As a consequence, answering cross-source
questions requires the system to reconcile different schemas and
granularities rather than querying a single homogeneous table.

All datasets used by Spark were curated into Parquet before the
Text-to-Spark-SQL experiments. The metadata layer therefore describes a
stable curated representation rather than exposing the language model
directly to raw source files.


\subsection{Yellow Taxi Trip Records}

The Yellow Taxi dataset is the largest source in the project and represents
the main data-volume component of the experimental setting. The January 2024
raw input contained 2,964,624 records. After curation, 2,963,713 records were
retained with 19 physical columns.

The dataset includes pickup and drop-off timestamps, pickup and drop-off
location identifiers, trip distance, fare-related attributes, and other
operational fields. In the metadata catalog, the physical Spark table is
registered as \texttt{yellow\_taxi}.

The pickup timestamp,
\texttt{tpep\_pickup\_datetime}, is particularly important for cross-source
analytics because it provides the temporal reference used to associate a
trip with the corresponding hourly weather observation. Similarly,
\texttt{PULocationID} and \texttt{DOLocationID} provide the physical keys
used to connect taxi trips with geographic information from the Taxi Zones
dataset.


\subsection{Green Taxi Trip Records}

The Green Taxi dataset provides a second trip-record source with a similar
analytical domain but a different physical schema. The January 2024 raw input
contained 56,551 records, of which 56,429 remained after curation. The
curated representation contains 20 physical columns and is registered in
Spark as \texttt{green\_taxi}.

Although Yellow and Green Taxi describe related events, they cannot be
treated as a single identical schema. In particular, their pickup and
drop-off timestamp columns use different physical names. Green Taxi uses
\texttt{lpep\_pickup\_datetime} for the pickup timestamp, whereas Yellow
Taxi uses \texttt{tpep\_pickup\_datetime}. This difference illustrates the
role of semantic metadata: both fields represent the same general concept
of a pickup time while remaining distinct physical SQL columns.

As for Yellow Taxi, Green Taxi location identifiers can be related to Taxi
Zones, while the pickup timestamp can be converted to an hourly bucket and
related to the weather source.


\subsection{Taxi Zones}

Taxi Zones is a compact geographic reference dataset containing 265 rows and
4 physical columns. It maps a numeric \texttt{LocationID} to descriptive
geographic attributes including zone and borough information.

Despite its small physical size, Taxi Zones is essential for many analytical
questions. Taxi-trip records store location identifiers rather than
human-readable geographic names. Therefore, a question such as ``Which
borough received the largest number of Green Taxi drop-offs?'' requires a
join between the transactional trip source and the zone reference source.

The project explicitly models four taxi-to-zone relationships: Yellow pickup,
Yellow drop-off, Green pickup, and Green drop-off. Distinguishing pickup from
drop-off relationships is necessary because both are represented by location
identifiers but correspond to different analytical meanings.


\subsection{Hourly Weather}

The weather source provides the temporal environmental dimension used for
questions involving rain, temperature, and hourly conditions. The curated
January 2024 dataset contains 744 hourly buckets and 23 columns.

Of the 744 expected hourly positions for the month, 743 contain an observed
weather record. One hourly bucket is explicitly represented as missing rather
than silently removed. The curated data contains 167 rainy hours and 576 dry
hours, with one hour whose precipitation status remains unknown because of
the missing observation.

Weather observations are registered in Spark as
\texttt{weather\_hourly}. The physical \texttt{weather\_hour} column is the
temporal key used for joins. Taxi pickup timestamps are converted to hourly
granularity using \texttt{date\_trunc('hour', ...)} before being compared
with this field.

Precipitation semantics are normalized during curation. Positive
precipitation values are treated as rainy conditions. Trace precipitation is
also considered precipitation even though no positive numeric amount can be
assigned to it. A zero precipitation amount represents a dry observation,
while missing precipitation information is kept distinct from both rainy and
dry conditions.


\subsection{Data Curation}

Curation was performed before metadata indexing and query generation so that
all downstream experiments operate over stable physical datasets.

For the taxi sources, records were restricted to pickups occurring in January
2024. Rows with missing pickup or drop-off timestamps were removed, and a
drop-off was required to occur after its corresponding pickup. Pickup and
drop-off location identifiers were required to correspond to valid taxi-zone
identifiers. Trip distance was constrained to non-negative values not
exceeding 1,000 miles.

The curation procedure intentionally avoids applying overly aggressive domain
assumptions. Zero-distance trips are preserved, as are negative monetary
values and zero passenger counts when they are present in the source data.
These values can represent unusual but potentially meaningful operational
records and were therefore not removed merely because they appear
counter-intuitive.

For weather data, raw observations were normalized to one record per clock
hour. When multiple observations were available within the same hourly
bucket, the preferred observation type was selected when available, with a
deterministic fallback to the latest suitable observation. The resulting
representation provides a stable hourly temporal key for cross-source joins.


\subsection{Cross-Source Relationships}

The metadata catalog explicitly represents six relationships required by the
core data lake.

Four relationships connect taxi locations to Taxi Zones:

\begin{itemize}

    \item \texttt{yellow\_taxi.PULocationID} to
    \texttt{taxi\_zones.LocationID};

    \item \texttt{yellow\_taxi.DOLocationID} to
    \texttt{taxi\_zones.LocationID};

    \item \texttt{green\_taxi.PULocationID} to
    \texttt{taxi\_zones.LocationID};

    \item \texttt{green\_taxi.DOLocationID} to
    \texttt{taxi\_zones.LocationID}.

\end{itemize}

Two additional relationships connect taxi pickup times with hourly weather:

\begin{itemize}

    \item
    \texttt{date\_trunc('hour',
    yellow\_taxi.tpep\_pickup\_datetime)}
    to \texttt{weather\_hourly.weather\_hour};

    \item
    \texttt{date\_trunc('hour',
    green\_taxi.lpep\_pickup\_datetime)}
    to \texttt{weather\_hourly.weather\_hour}.

\end{itemize}

These relationships are stored as metadata rather than being inferred by the
language model for every request. This allows relation-aware retrieval to
expand an initially selected dataset or column into the physical join
endpoints required for execution.


\subsection{Volume and Variety in the Experimental Data}

The data lake represents the two Big Data dimensions selected for this
project.

\emph{Volume} is primarily provided by the taxi-trip datasets and is evaluated
experimentally using progressively larger subsets of Yellow Taxi data, from
100,000 records to the complete curated set of 2,963,713 records.

\emph{Variety} arises from the combination of transactional mobility data,
geographic reference data, and hourly weather observations. The sources have
different schemas, column naming conventions, row granularities, and
relationships. This heterogeneity motivates the metadata-aware architecture:
a natural-language request can refer to one conceptual analytical task while
its execution requires physical fields distributed across multiple datasets.
